\documentclass[a4paper,12pt]{article}

\usepackage[a4paper,margin=1in]{geometry}
\usepackage{subcaption}

\setlength{\parskip}{12pt}
\setlength{\parindent}{0pt}

\usepackage[utf8]{inputenc}

\usepackage{amsmath}
\usepackage[inkscapearea=drawing]{svg}
\usepackage{amsthm}
\usepackage{amsfonts}
\usepackage{amssymb}
\usepackage{mathtools}
\usepackage{hyperref}
\usepackage{needspace}
\usepackage{letltxmacro}
\usepackage{graphicx}
\usepackage{extarrows}

% easily make a set 
\DeclarePairedDelimiterX\set[1]{\lbrace}{\rbrace}{\def\given{\;\delimsize\vert\;}\,#1\,}

% index concatenation
\newcommand\doubleplus{\mathbin{+\!+}}

% wrap single quotes around a symbol in math mode, used for index parts
\newcommand{\HSindexpart}[1]{\text{`}#1\text{'}}

% used commonly

\newcommand{\contradiction}{\Rightarrow\!\Leftarrow}

\newcommand{\HSview}[1][V]{\widetilde{#1}}

\newcommand{\HStriangle}[1][I]{\triangle_{\HSview:#1}}

\newcommand{\HSA}[1][I]{A_{\HSview:#1}}
\newcommand{\HSB}[1][I]{B_{\HSview:#1}}
\newcommand{\HSC}[1][I]{C_{\HSview:#1}}
\newcommand{\HSAx}[1][I]{A_{x\HSview:#1}}
\newcommand{\HSBx}[1][I]{B_{x\HSview:#1}}
\newcommand{\HSCx}[1][I]{C_{x\HSview:#1}}
\newcommand{\HSAy}[1][I]{A_{y\HSview:#1}}
\newcommand{\HSBy}[1][I]{B_{y\HSview:#1}}
\newcommand{\HSCy}[1][I]{C_{y\HSview:#1}}

\newcommand{\HST}[1][I]{T_{\HSview:#1}}
\newcommand{\HSL}[1][I]{L_{\HSview:#1}}
\newcommand{\HSR}[1][I]{R_{\HSview:#1}}
\newcommand{\HSTx}[1][I]{T_{x\HSview:#1}}
\newcommand{\HSLx}[1][I]{L_{x\HSview:#1}}
\newcommand{\HSRx}[1][I]{R_{x\HSview:#1}}
\newcommand{\HSTy}[1][I]{T_{y\HSview:#1}}
\newcommand{\HSLy}[1][I]{L_{y\HSview:#1}}
\newcommand{\HSRy}[1][I]{R_{y\HSview:#1}}

\newcommand{\HSPSI}[1][I]{\Psi_{\HSview:#1}}
\newcommand{\HSPSIx}[1][I]{\Psi_{x\HSview:#1}}
\newcommand{\HSPSIy}[1][I]{\Psi_{y\HSview:#1}}

\newcommand{\HSPHI}[1][I]{\Phi_{\HSview:#1}}
\newcommand{\HSPHIx}[1][I]{\Phi_{x\HSview:#1}}
\newcommand{\HSPHIy}[1][I]{\Phi_{y\HSview:#1}}

\newcommand{\HSO}[1][I]{O_{\HSview:#1}}
\newcommand{\HSOx}[1][I]{O_{x\HSview:#1}}
\newcommand{\HSOy}[1][I]{O_{y\HSview:#1}}

\newcommand{\HStau}{\tau_{\HSview}}
\newcommand{\HSlambda}{\lambda_{\HSview}}
\newcommand{\HSrho}{\rho_{\HSview}}

\newcommand{\HSk}{k_{\HSview}}
\newcommand{\HSdelta}[1][d]{\delta_{#1}}
\newcommand{\HSs}[1][d]{s_{\HSview:#1}}
\newcommand{\HSsigma}[1][d]{\sigma_{\HSview:#1}}
\newcommand{\HSkappa}[1][d]{\kappa_{\HSview:#1}}


% insert a thin space after a comma in math mode, if it is followed by a space in the source
% stolen from stack exchange
\makeatletter
\AtBeginDocument{%
  \mathchardef\stdcomma=\mathcode`,
  \mathcode`,="8000
}
\begingroup\lccode`~=`, \lowercase{\endgroup\def~}{%
  \futurelet\@let@token\spaced@comma
}
\newcommand{\spaced@comma}{%
  \stdcomma
  \ifx\@let@token\@sptoken\,\fi
}
\makeatother

% closed square root, i slanted it a bit because i think it looks pretty
% also stolen from stack exchange and modified slightly
\makeatletter
\let\oldr@@t\r@@t
\def\r@@t#1#2{%
\setbox0=\hbox{$\oldr@@t#1{#2}$}\dimen0=\ht0
\advance\dimen0-0.2\ht0
\setbox2=\hbox{\kern-0.025em\rotatebox[origin=t]{20}{\vrule height\ht0 depth -\dimen0}}%
{\box0\lower0.6pt\box2}}
\LetLtxMacro{\oldsqrt}{\sqrt}
\renewcommand*{\sqrt}[2][\ ]{\oldsqrt[#1]{#2} }
\makeatother

\newtheorem{theorem}{Theorem}[section]
\newtheorem{lemma}[theorem]{Lemma}
\newtheorem{proposition}[theorem]{Proposition}
\newtheorem*{remark}{Remark}
\newtheorem{conjecture}[theorem]{Conjecture}

\theoremstyle{definition}
\newtheorem{definition}[theorem]{Definition}
\theoremstyle{plain}

\title{Hell Sangaku}
\author{NyteLyte}
\date{}

\begin{document}
	\maketitle

	\tableofcontents
	\pagebreak

	\section{Introduction}

	Starting with a circle around the origin, inscribe an arbitrary non-degenerate triangle into it.
	Our goal is to put more copies of that triangle (where only scaling and translation are
	allowed) into the circle such that the following must be satisfied at every step:\begin{itemize}
		\item every triangle must touch the circle at least once,
		\item every triangle must touch either the circle or another triangle at all three of its vertices,
		\item the triangles must not overlap (neither interiors nor edges) and they must not share vertices
	\end{itemize}

	We put as many copies of the triangle into the circle as we can.

	\begin{figure}[h]
		\centering
		\def\svgwidth{\columnwidth}
		\scalebox{0.85}{\includesvg{figure1.svg}}
		\caption{The fractal this procedure creates, using an equilateral triangle}
	\end{figure}

	\newpage
	\section{Definitions and notation}

	A point $P$ is an ordered pair of 2 real numbers $x$ and $y$.
	Referring to a specific coordinate on a point is written $P_x$ and $P_y$.
	
	We write the set of all triangles as $\mathbb{T} \coloneqq \set{ (A, B, C) \given A, B, C \in \mathbb{R}^2 }$.

	We write the set of all triangles where the winding order is counter-clockwise as $\mathbb{T}^+$. All triangles from now on will be assumed to be in $\mathbb{T}^+$.

	We also write some triangle $\left(A, B, C\right)$ as $\triangle ABC$.
	
	\begin{definition}
		A \textbf{view transformation} $\widetilde{P}_{\triangle ABC} : \mathbb{R}^2 \to \mathbb{R}^2$, where $\triangle ABC$ is a triangle and $P \in \set{A, B, C}$,
		is a rotation of the coordinate space around the origin such that $P_y > Q_y = R_y$, where $Q$ and $R$ are the two triangle points that are $\neq P$.
	\end{definition}
	
	\begin{definition}
		An \textbf{inverse view transformation} $\widetilde{P}_{\triangle ABC}^{-1} : \mathbb{R}^2 \to \mathbb{R}^2$
		of view transformation $\widetilde{P}_{\triangle ABC}$ is a function such that
		$$\widetilde{P}_{\triangle ABC}^{-1}\left(\widetilde{P}_{\triangle ABC}\left(p\right)\right) =
		  \widetilde{P}_{\triangle ABC}\left(\widetilde{P}_{\triangle ABC}^{-1}\left(p\right)\right) =
		  p, \forall p \in \mathbb{R}^2$$
	\end{definition}

	A view transformation applied to a triangle is a shorthand for a triangle with the view transformation applied to all 3 of its points.
	For all view transformations $\HSview$, $$\HSview\left(\triangle ABC\right) \coloneqq \left(\HSview\left(A\right), \HSview\left(B\right), \HSview\left(C\right)\right)$$
	This also applies to inverse view transformations.

	The radius of the circle is $r$.

	\begin{definition}
		The \textbf{root triangle} $\triangle_R \coloneqq \triangle A_R B_R C_R$ is a triangle touching the circle at all 3 points.
		It has angles $\alpha$, $\beta$, and $\gamma$ at vertices $A_R$, $B_R$, and $C_R$, respectively.
	\end{definition}

	\begin{conjecture}\label{uniqueness-conjecture}
		For any choice of a root triangle---that is, a choice of 3 distinct points on the circle---there is a unique solution for the set of triangles that satisfies the constraints of the problem.
		It is possible to fit an infinite amount of triangles into the circle this way.
	\end{conjecture}

	We define the three canonical view transformations in relation to the root triangle:
	$$\widetilde{A} \coloneqq {\widetilde{A_R}} _ {\triangle_R}$$
	$$\widetilde{B} \coloneqq {\widetilde{B_R}} _ {\triangle_R}$$
	$$\widetilde{C} \coloneqq {\widetilde{C_R}} _ {\triangle_R}$$

	\begin{remark}
		Because of the constraints of the problem---all triangles filling the circle must be of the same rotation---those three
		view transformations are the only relevant ones in the problem, and thus may apply to any subsequent triangle.
	\end{remark}

	After a view transformation has been applied to a triangle, we will
	refer to such a triangle as a \textbf{positive triangle}.
	When we speak of a positive triangle, we refer to its points as ``top'', ``left'', ``right'', defined later.

	\begin{remark}
		The triangle is referred to as positive because it has a clear direction it is ``pointing'' in, upwards.
		All a positive triangle really is is a regular triangle in a view transformation with vertices ``top'', ``left'', ``right''.
	\end{remark}

	\begin{figure}[h]
		\begin{subfigure}{0.5\textwidth}
			\centering
			\def\svgwidth{\columnwidth}
			\scalebox{0.85}{\includesvg{figure2a.svg}}
			\caption{Before applying $\HSview[C]$}
		\end{subfigure}
		\begin{subfigure}{0.5\textwidth}
			\centering
			\def\svgwidth{\columnwidth}
			\scalebox{0.85}{\includesvg{figure2b.svg}}
			\caption{After applying $\HSview[C]$}
		\end{subfigure}
		\caption{Application of a view transformation}
	\end{figure}
	
	\begin{remark}
		Notice that, after inscribing the root triangle into the circle, we have essentially `split' the circle into a triangle and three circular segments.
		A view transformation of one of the root triangle's points places the segment opposite that point below the root triangle in that transformed space.
	\end{remark}
	
	\begin{definition}
		A triangle $\triangle_X$ is said to \textbf{touch} triangle $\triangle_Y$ if a vertex of $\triangle_X$ touches a side of $\triangle_Y$.
		A triangle is said to touch the circle if one of its vertices is on the circle.
	\end{definition}
	
	\begin{remark}
		Triangle $\triangle_X$ touching triangle $\triangle_Y$ does not imply that $\triangle_Y$ touches $\triangle_X$.
		In fact, since triangles cannot share vertices or overlap in our problem, it implies the opposite:
		if $\triangle_X$ touches $\triangle_Y$, $\triangle_Y$ cannot touch $\triangle_X$.
	\end{remark}
	
	After applying some view transformation $\HSview$, we have to construct a triangle that, at all 3 of its points, touches either the circle or another triangle,
	with at least one of those being the circle. It is clear that, since the bottom side is (by definition of $\HSview$) parallel to the $x$ axis, and it is below the
	root triangle, both of the bottom points of our constructed triangle must touch the circle.
	\begin{conjecture}
		The top point of our constructed triangle must touch the triangle above it and cannot touch the circle.
	\end{conjecture}

	\begin{definition}
		A \textbf{base triangle} is a triangle inscribed in a circular segment such that two of its vertices
		touch circle and its remaining vertex touches another base triangle or the root triangle.
	\end{definition}
	
	\begin{figure}[h]
		\centering
		\def\svgwidth{\columnwidth}
		\scalebox{0.65}{\includesvg{figure3.svg}}
		\caption{The base triangle inscribed in a circular segment created by the root triangle}
	\end{figure}

	\begin{remark}
		Notice that, after inscribing the base triangle into the circular segment, we have essentially `split' the circular segment into a triangle, another circular segment, and two other regions.
		If we were to inscribe another triangle in the segment, we would get another base triangle.
	\end{remark}

	\begin{definition}
		A \textbf{pseudo-sector} is a region bounded by a circular arc and two line segments from the arc's endpoints to an arbitrary interior point.
		The arbitrary interior point is called the \textbf{origin} of the pseudo-sector.
	\end{definition}

	\begin{remark}
		A pseudo-sector in our problem is the space enclosed by an arc and 2 triangles, where the ``arbitrary interior point'' is the point where one triange touches another.
		Its non-arc components, thus, lie on the sides of 2 different triangles. Since we are in a view transformation, one of the line segments is always on top and parallel to the $x$ axis.
		We can call that line segment the top part of the pseudo-sector. The other line segment will be the left/right part of the pseudo-sector.
	\end{remark}

	We now have 2 pseudo-sectors to choose from, the left pseudo-sector and right pseudo-sector.

	Notice how, if we were to extend a horizontal line from the circular arc belonging to a pseudosector toward the
	inside of the circle, it would intsersect another edge of the pseudo-sector before intersecting
	the other side of the circle, because the pseudo-sector covers the entire vertical area we are drawing the horizontal line from.
	This applies both to the left and right pseudo-sectors.

	Therefore, if we were to construct a triangle inside a pseudo-sector, its bottom side would have to touch the circle at least once, and another triangle at least once.

	\begin{conjecture}
		The top point of our constructed triangle must touch the top part of the pseudo-sector, which is a bottom line of another triangle, different from the one it already has to touch.
		It cannot touch the circle.
	\end{conjecture}

	\begin{definition}
		A \textbf{normal triangle} is a triangle inscribed in a pseudo-sector such that it touches the circle at one vertex, and two different triangles at its other two vertices.
	\end{definition}
	
	\begin{figure}[h]
		\centering
		\def\svgwidth{\columnwidth}
		\scalebox{0.65}{\includesvg{figure4.svg}}
		\caption{A normal triangle inscribed in a left pseudo-sector of a base triangle}
	\end{figure}
	
	\begin{remark}
		Notice that, after inscribing the normal triangle into the pseudo-sector, we have essentially `split' the pseudo-sector into two triangles and two pseudo-sectors.
		If we were to inscribe another triangle in either of the two pseudo-sectors, we would get another normal triangle.
		Notice how one of the pseudo-sectors is below the inscribed triangle, and one is to the left/right (horizontal), depending on which main pseudo-sector we chose from the base.
	\end{remark}

	\begin{definition}
		A \textbf{horizontal pseudo-sector} of a normal triangle is the pseudo-sector to its left, on the left side of the circle, or the pseudo-sector to its right, on the right side of the circle.
	\end{definition}

	\begin{definition}
		A \textbf{vertical pseudo-sector} of a normal triangle is the pseudo-sector below it.
	\end{definition}

	\begin{definition}
		A \textbf{horizontal triangle} is a normal triangle constructed inside a horizontal pseudo-sector of a normal or a left/right pseudo-sector of a base triangle.
	\end{definition}

	\begin{definition}
		A \textbf{vertical triangle} is a normal triangle constructed inside a vertical pseudo-sector of a normal triangle.
	\end{definition}

	\Needspace{15\baselineskip}
	Notice how, at each step, we are creating new spaces to fit more triangles into. As such, we can establish a parent-child relationship and structure this problem as a tree.
	\begin{definition}
		A \textbf{triangle tree} is a tree with the following node types:\begin{itemize}
			\item a \textbf{root triangle node} which represents the root triangle,
			\item a \textbf{root side node} which represents a side of the root triangle (a/b/c),
			\item a \textbf{base triangle node} which represents a base triangle,
			\item a \textbf{base side node} which represents a side of the base triangle (left/right),
			\item a \textbf{normal triangle node} which represents a normal triangle
		\end{itemize}
		
		A root triangle node has no parents and has three root side node children.
		These children are called `a', `b', and `c'.

		A root side node has a root triangle node parent and one base triangle node child.
		This child is called `base'.

		A base triangle node has either a root side node parent or a base triangle node parent, two base side node children, and one base triangle node child.
		The two base side node children are called `left' and `right', and the base triangle node child is called `base'

		A base side node has a base triangle node parent and one normal triangle node child.
		This child is called `horizontal'.

		A normal triangle node has either a base side node parent or a normal triangle node parent, and two normal triangle node children.
		These children are called `horizontal' and `vertical'.

	\end{definition}

	\Needspace{15\baselineskip}
	We also need a way to refer to a given triangle.
	\begin{definition}
		An \textbf{index part} represents a triangle tree node type and is one of the following symbols: \begin{itemize}
			\item `$R$', which represents a root triangle node
			\item `$a$'/`$b$'/`$c$', which represents a root side node
			\item `$B$', which represents a base triangle node
			\item `$l$'/`$r$', which represents a base side node, where `$l$' is left and `$r$' is right
			\item `$H$'/`$V$', which represents a normal triangle node, where `$H$' is horizontal and `$V$' is vertical
		\end{itemize}
	\end{definition}

	\Needspace{15\baselineskip}
	\begin{definition}
		An \textbf{index} is a string of concatenated string of identifier parts, representing a path taken in a triangle tree.
		It is subject to the following restrictions:\begin{itemize}
			\item it must be nonempty,
			\item an `$R$' part must be preceeded by nothing,
			\item an `$a$'/`$b$'/`$c$' part must be preceeded by an `$R$' part,
			\item a `$B$' part must be preceeded by an `$a$'/`$b$'/`$c$' part or a `$B$' part,
			\item an `$l$'/`$r$' part must be preceeded by a `$B$' part,
			\item an `$H$' part must be preceeded by an `$l$'/`$r$' part, an `$H$' part, or a `$V$' part,
			\item a `$V$' part must be preceeded by an `$H$' part or a `$V$' part
		\end{itemize}

		Concatenation of an index part to an index can also be denoted by $I \doubleplus i$, where $I$ is an index and $i$ is an index part.
		This will be used where juxtaposition is ambiguous, like with a variable and a symbol (e.g. $I \doubleplus \HSindexpart{H}$ is preferred to $IH$).

		An index is said to be real if it ends with an `$R$', `$B$', or an `$H$'/`$V$' part.
	\end{definition}

	\begin{definition}
		The \textbf{root index} is index $R$.
		A \textbf{base index} is an index which ends with a `$B$' part.
		A \textbf{normal index} is an index which ends with a `$H$'/`$V$' part.
		A \textbf{horizontal index} is an index which ends with a `$H$' part.
		A \textbf{vertical index} is an index which ends with a `$V$' part.
		A \textbf{left index} is an index which contains an `$l$' part.
		A \textbf{right index} is an index which contains an `$r$' part.
	\end{definition}

	Examples of indices are:\\
	$R$: the root triangle,\\
	$RaB$: the first base triangle on root side a,\\
	$RbB$: the first base triangle on root side b,\\
	$RaBBlHHV$: the vertical child of the horizontal child of the left horizontal child of the base child of the base child of the root triangle on root side a
	
	With $I$ being a real index and $\HSview$ being a view transformation:\\
	$\triangle_I \coloneqq \triangle A_I B_I C_I$: a triangle in the triangle tree without any transformations,\\
	$\HSview\left(\triangle_I\right) \coloneqq \HStriangle \coloneqq \triangle \HSA \HSB \HSC$: a triangle in the triangle tree with a view transformation

	We also denote the following substitutions:

	$\HST = \begin{cases}
		\HSA, & \HSview \text{ is } \HSview[A]\\
		\HSB, & \HSview \text{ is } \HSview[B]\\
		\HSC, & \HSview \text{ is } \HSview[C]\\
  \end{cases}$,\quad\quad
	$\HSL = \begin{cases}
		\HSB, & \HSview \text{ is } \HSview[A]\\
		\HSC, & \HSview \text{ is } \HSview[B]\\
		\HSA, & \HSview \text{ is } \HSview[C]\\
  \end{cases}$,\quad\quad
	$\HSR = \begin{cases}
		\HSC, & \HSview \text{ is } \HSview[A]\\
		\HSA, & \HSview \text{ is } \HSview[B]\\
		\HSB, & \HSview \text{ is } \HSview[C]\\
  \end{cases}$
	
	\pagebreak
	We analogously define the following angle substitutions:

	$\tau_{\HSview} = \begin{cases}
		\alpha, & \HSview \text{ is } \HSview[A]\\
		\beta, & \HSview \text{ is } \HSview[B]\\
		\gamma, & \HSview \text{ is } \HSview[C]\\
  \end{cases}$,\quad\quad
	$\lambda_{\HSview} = \begin{cases}
		\beta, & \HSview \text{ is } \HSview[A]\\
		\gamma, & \HSview \text{ is } \HSview[B]\\
		\alpha, & \HSview \text{ is } \HSview[C]\\
  \end{cases}$,\quad\quad
	$\rho_{\HSview} = \begin{cases}
		\gamma, & \HSview \text{ is } \HSview[A]\\
		\alpha, & \HSview \text{ is } \HSview[B]\\
		\beta, & \HSview \text{ is } \HSview[C]\\
  \end{cases}$
	
	If we are referring to a coordinate of some point on any triangle in the triangle tree, we write the
	coordinate before the index or view transformation in the subscript: $A_{xI}$, $\HSAx$, $\HSTx$.
	
	\begin{figure}[h]
		\centering
		\def\svgwidth{\columnwidth}
		\scalebox{0.65}{\includesvg{figure5.svg}}
		\caption{Triangle tree structure}
	\end{figure}
	
	We denote:\\
	$\mathcal{I}_\Omega$: the set of all indices,\\
	$\mathcal{I}$: the set of all real indices,\\
	$\mathcal{I}_R$: the set containing the root index,\\
	$\mathcal{I}_{B}$: the set containing every base index,\\
	$\mathcal{I}_{Bv}$: the set containing every base index on root side $v$,\\
	$\mathcal{I}_{N}$: the set containing every normal index,\\
	$\mathcal{I}_{Nvd}$: the set containing every normal index on root side $v$ and base side $d$\\
	$\mathcal{P}$: the set containing every index part,\\
	$\mathcal{P}_{R}$: the set containing the root index part\\
	$\mathcal{P}_{s}$: the set containing the root side index parts, `$a$', `$b$', and `$c$',\\
	$\mathcal{P}_{B}$: the set containing the base index part\\
	$\mathcal{P}_{d}$: the set containing the base side index parts, `$l$' and `$r$',\\
	$\mathcal{P}_{N}$: the set containing the normal index parts, `$H$' and `$V$'
	
	\begin{definition}
		The \textbf{parent function} $p : \mathcal{I}_\Omega \setminus \mathcal{I}_R \to \mathcal{I}_\Omega$ is defined by
		$$p(I \doubleplus i) = I$$
		where $I \in \mathcal{I}_\Omega, i \in \mathcal{P}$
	\end{definition}
	
	\begin{definition}
		The \textbf{real parent function} $P : \mathcal{I} \setminus \mathcal{I}_R \to \mathcal{I}$ is defined by
		$$P(I \doubleplus i) = \begin{cases}
			I, & I \in \mathcal{I}\\
			p(I), & \text{otherwise}
  	\end{cases}$$
		where $I \in \mathcal{I}_\Omega, i \in \mathcal{P}$
	\end{definition}
	
	It is useful to derive which triangles an arbitrary triangle touches.
	
	Suppose we are in some view transformed space $\HSview$.

	The root triangle, by definition, does not touch any other triangles.

	It is clear that \textbf{a base triangle always touches its parent at its top vertex}. Its top vertex is also the origin of the triangle's left and right pseudo-sectors.

	Notice now that, if we were to inscribe a horizontal triangle into one of the base triangle's pseudo-sectors,
	its top point would have to be on the base triangle's respective pseudo-sector's top side.

	Suppose we are on the left pseudo-sector side, then the left point of the constructed normal triangle touches the circle,
	its top point touches the pseudo-sector's top side, and its right point touches its parent base triangle.
	Notice how the top point is the origin of the horizontal pseudo-sector of the constructed normal triangle, and the right point is the origin of the vertical pseudo-sector.
	The argument is analogous for the right pseudo-sector side, left and right are just swapped.

	Since the origin of the normal triangle's horizontal pseudo-sector is on the same line as the origin of the pseudo-sector the normal triangle is contained in,
	we can apply the same argument recursively, and conclude that
	\textbf{a horizontal triangle's top vertex touches the same triangle that its parent's top vertex touches}.
	
	It is clear that \textbf{the other triangle that a horizontal triangle touches is its parent}.

	Since the origin of the normal triangle's vertical pseudo-sector is on the same line as the origin of the pseudo-sector the normal triangle is contained in, we can
	apply the same argument recursively, and conclude that
	\textbf{a vertical triangle's non-top non-circle vertex touches the same triangle that its parent's non-top non-circle vertex touches}.
	
	It is clear that \textbf{the other triangle that a vertical triangle touches is its parent}.

	\begin{definition}
		A \textbf{horizontal touching index function} $t_H : \mathcal{I}_N \to \mathcal{I}_N \cup \mathcal{I}_B$ is defined by
		$$t_H(I) = \begin{cases}
			P(I),      & I \text{ is a horizontal index}\\
			t_H(P(I)), & I \text{ is a vertical index}
		\end{cases}$$
	\end{definition}
	
	\begin{definition}
		A \textbf{vertical touching index function} $t_V : \mathcal{I}_N \cup \mathcal{I}_B \to \mathcal{I}$ is defined by
		$$t_V(I) = \begin{cases}
			P(I),      & I \text{ is a vertical index or a base index}\\
			t_V(P(I)), & I \text{ is a horizontal index}
		\end{cases}$$
	\end{definition}
	
	\begin{definition}
		For some view transformed space $\HSview$ and normal index $I$, we will denote the pseudo-sector origin point
		$$\HSO = \begin{cases}
			\HST[P(I)], & I \text{ is a horizontal index}\\
			\HSL[P(I)], & I \text{ is a right index and a vertical index}\\
			\HSR[P(I)], & I \text{ is a left index and a vertical index}
		\end{cases}$$
		The pseudo-sector origin point of some normal triangle is the origin point of the pseudo-sector that normal
		triangle is inscribed in.
	\end{definition}

	\begin{remark}
		Notice how the $y$ coordinate of the top point of any normal triangle with index $I$ is its pseudo-sector's origin point's $y$ coordinate.
		Notice, how, on the left base side, the normal triangle with index $I$ touches the line of angle $\HSlambda$, which goes through $\HSO$.
		Notice, how, on the right base side, the normal triangle with index $I$ touches the line of angle $-\HSrho$, which goes through $\HSO$.
		Notice how, this point is always inside the circle and not on it.
	\end{remark}
	

	Every triangle has a side-to-height ratio on each of its sides. This side-to-height ratio remains constant for similar triangles.
	In some view transformed space $\HSview$, the side-to-height ratio relevant to the problem is always going to be the
	one relating to the triangle's bottom side, so we write
	$$\HSk \coloneqq \cot\left(\HSlambda\right) + \cot\left(\HSrho\right)$$

	Additionally, for some view transformation $\HSview$ and base side index part $d \in \mathcal{P}_d$
	we will define the following constant: the slope of the left/right side when represented as a function of $y$:
	$$\HSs \coloneqq \begin{cases}\cot\left(\HSlambda\right),& d = \HSindexpart{l}\\-\cot\left(\HSrho\right),& d = \HSindexpart{r}\end{cases}$$
	
	\newpage

	\section{The root triangle}
	We are given 2 angles $\alpha$ and $\beta$, corresponding to the angles of a triangle at vertices $A$ and $B$.
	The angle at vertex $C$, $\gamma$, is redundent, so it does not need to be given.
	We are also given an angle $\theta$, which is an offset rotation angle.

	We place the initial point $A$ at $\left(r\cos\left(\theta\right), r\sin\left(\theta\right)\right)$.

	In order to get the angles of the triangle to be $\alpha$, $\beta$, and $\gamma$, we use the fact that the angle
	at the center of the circle is twice the angle at the circumferance. Our angles at the circumferance are
	$\alpha$, $\beta$, and $\gamma$, respectively.

	Since we are requiring that our circle have counter-clockwise winding order, our vertices have to be counter-clockwise
	in order $A$, $B$, $C$. We know that the angle opposite side $c$, connecting $A$ and $B$, has to be $\gamma$.
	Therefore, the arc distance between $A$ and $B$ has to be $2\gamma$. Therefore,
	$B = \left(r\cos\left(\theta + 2\gamma\right), r\sin\left(\theta + 2\gamma\right)\right)$.
	We know that the angle opposite $a$, connecting $B$ and $C$, has to be $\alpha$.
	Therefore, the arc distance between $B$ and $C$ has to be $2\alpha$.
	Therefore,
	$C = \left(r\cos\left(\theta + 2\gamma + 2\alpha\right), r\sin\left(\theta + 2\gamma + 2\alpha\right)\right)$.
	Note that $\alpha + \beta + \gamma = \pi$.
	Therefore, $2\gamma + 2\alpha = 2\pi - 2\beta$. Therefore,
	$C = \left(r\cos\left(\theta + 2\pi - 2\beta\right), r\sin\left(\theta + 2\pi - 2\beta\right)\right)$. Because
	$\sin$ and $\cos$ both have a period of $2\pi$, we can further simplify this to:
	$C = \left(r\cos\left(\theta - 2\beta\right), r\sin\left(\theta - 2\beta\right)\right)$.

	We now have the following points:
	\begin{align*}
		A =& \left(r\cos\left(\theta\right), r\sin\left(\theta\right)\right) \\
		B =& \left(r\cos\left(\theta + 2\gamma\right), r\sin\left(\theta + 2\gamma\right)\right) \\
		C =& \left(r\cos\left(\theta - 2\beta\right), r\sin\left(\theta - 2\beta\right)\right)
	\end{align*}

	Now that we have the root triangle that is not in any view, we need to construct the view transformations
	$\HSview[A]$, $\HSview[B]$, $\HSview[C]$.

	Since all of the points are on the circle centered around the origin, and their coordinates are already in the form
	$\left(r\cos\left(\Theta_1\right), r\sin\left(\Theta_1\right)\right)$, their rotation by some angle $\Theta_2$ is
	equivalent to simply adding the angles:
	$\left(r\cos\left(\Theta_1 + \Theta_2 \right), r\sin\left(\Theta_1 + \Theta_2\right)\right)$.

	Starting with view $\HSview[C]$, we require that $c$ is parallel to the $x$-axis and that $C$ is above $c$ on the
	$y$-axis. Since a view is a rotation of the coordinate system, we will define the angle by which we need to rotate
	the coordinate system with $\theta_{\HSview[C]}$.

	Without loss of generality, let $\theta = 0$. If $\theta \neq 0$, we can
	simply add $-\theta$ to $\theta_{\HSview[C]}$ to undo its rotation.

	Let $A' = \HSview[C]\left(A\right), B' = \HSview[C]\left(B\right), C' = \HSview[C]\left(C\right)$.

	We require the following:
	$$
		C'_y > A'_y = B'_y \iff
		r\sin\left(2\gamma + 2\alpha + \theta_{\HSview[C]}\right) >
		r\sin\left(\theta_{\HSview[C]}\right) =
		r\sin\left(2\gamma + \theta_{\HSview[C]}\right)
	$$

	If we let $\theta_{\HSview[C]} = \alpha + \beta + \frac{\pi}{2}$, we get
	\begin{align*}
		& r\sin\left(\alpha + \beta + \frac{\pi}{2}\right) = 
		  r\sin\left(2\gamma + \alpha + \beta + \frac{\pi}{2}\right) & \iff \\
		& \sin\left(\pi - \gamma + \frac{\pi}{2}\right) = 
		  \sin\left(\gamma + \pi + \frac{\pi}{2}\right) & \iff \\
		& \sin\left(\frac{3\pi}{2} - \gamma\right) = 
		  \sin\left(\frac{3\pi}{2} + \gamma\right) & \iff \\
		& \sin\left(-\frac{\pi}{2} - \gamma\right) = 
		  \sin\left(\gamma - \frac{\pi}{2}\right) & \iff \\
		& -\sin\left(\frac{\pi}{2} + \gamma\right) = 
		  \sin\left(\gamma - \frac{\pi}{2}\right) & \iff \\
		& \sin\left(\frac{\pi}{2} + \gamma - \pi\right) = 
		  \sin\left(\gamma - \frac{\pi}{2}\right)
	\end{align*}

	Now, we need to prove that $C'_y > A'_y$ or $C'_y > B'_y$.

	\begin{align*}
		& r\sin\left(2\gamma + 2\alpha + \alpha + \beta + \frac{\pi}{2}\right) >
		  r\sin\left(\alpha + \beta + \frac{\pi}{2}\right) & \iff \\
		& r\sin\left(2\pi - 2\beta + \alpha + \beta + \frac{\pi}{2}\right) >
		  r\sin\left(\alpha + \beta + \frac{\pi}{2}\right) & \iff \\
		& r\sin\left(2\pi - \beta + \alpha + \frac{\pi}{2}\right) >
		  r\sin\left(\alpha + \beta + \frac{\pi}{2}\right) & \iff \\
		& r\sin\left(- \beta + \alpha + \frac{\pi}{2}\right) >
		  r\sin\left(\alpha + \beta + \frac{\pi}{2}\right) & \iff \\
		& r\sin\left(\alpha + \frac{\pi}{2} - \beta\right) >
		  r\sin\left(\alpha + \frac{\pi}{2} + \beta\right)
	\end{align*}

	It is sufficient to prove that, for all $\alpha, \beta \in \left(0, \pi\right)$, 
	$\sin\left(\alpha + \frac{\pi}{2} - \beta\right) >
	 \sin\left(\alpha + \frac{\pi}{2} + \beta\right)$.
	Rearranging, we get
	$\sin\left(\alpha + \frac{\pi}{2} - \beta\right) - \sin\left(\alpha + \frac{\pi}{2} + \beta\right) > 0$.

	We can use the product-sum identity for sin to get the following:
	\begin{align*}
		& 2\cdot\cos\left(\frac{\alpha + \frac{\pi}{2} - \beta + \alpha + \frac{\pi}{2} + \beta}{2}\right) \cdot
	    \sin\left(\frac{\alpha + \frac{\pi}{2} - \beta - \alpha - \frac{\pi}{2} - \beta}{2}\right) > 0 & \iff \\
		& \cos\left(\frac{2\alpha + \pi}{2}\right) \cdot
		  \sin\left(-\beta\right) > 0 & \iff \\
		& -\cos\left(\frac{2\alpha + \pi}{2}\right) \cdot
		  \sin\left(\beta\right) > 0 & \iff \\
	  & \cos\left(\frac{2\alpha + \pi}{2}\right) \cdot
	    \sin\left(\beta\right) < 0
	\end{align*}


	For $\beta \in \left(0, \pi\right)$, $\sin\left(\beta\right) > 0$, so we need to prove
	$\cos\left(\alpha + \frac{\pi}{2}\right) < 0$.

	Since 
	$\alpha \in \left(0, \pi\right)$, $\alpha + \frac{\pi}{2} \in \left(\frac{\pi}{2}, \frac{3\pi}{2}\right)$.
	Cosine is negative in that domain.

	Therefore, we have $\theta_{\HSview[C]} = \alpha + \beta + \frac{\pi}{2} - \theta = \pi - \gamma + \frac{\pi}{2} - \theta = \frac{3\pi}{2} - \gamma - \theta$.

	The next view we observe is $\HSview[B]$, we require that $b$ is parallel to the $x$-axis and that $B$ is above $b$ on the
	$y$-axis. Since a view is a rotation of the coordinate system, we will define the angle by which we need to rotate
	the coordinate system with $\theta_{\HSview[B]}$.
	
	Without loss of generality, let $\theta = 0$. If $\theta \neq 0$, we can
	simply add $-\theta$ to $\theta_{\HSview[B]}$ to undo its rotation.

	Let $A' = \HSview[B]\left(A\right), B' = \HSview[B]\left(B\right), C' = \HSview[B]\left(C\right)$.

	We require the following:
	$$
		B'_y > A'_y = C'_y \iff
		r\sin\left(2\gamma + \theta_{\HSview[B]}\right) >
		r\sin\left(\theta_{\HSview[B]}\right) =
		r\sin\left(2\gamma + 2\alpha + \theta_{\HSview[B]}\right)
	$$
	
	If we let $\theta_{\HSview[B]} = -\alpha - \gamma + \frac{\pi}{2}$, we get
	\begin{align*}
		& r\sin\left(\theta_{\HSview[B]}\right) =
		  r\sin\left(2\gamma + 2\alpha + \theta_{\HSview[B]}\right) & \iff \\
		& r\sin\left(-\alpha - \gamma + \frac{\pi}{2}\right) =
		  r\sin\left(2\gamma + 2\alpha + -\alpha - \gamma + \frac{\pi}{2}\right) & \iff \\
		& r\sin\left(-\alpha - \gamma + \frac{\pi}{2}\right) =
		  r\sin\left(\gamma + \alpha + \frac{\pi}{2}\right) & \iff \\
		& r\sin\left(\frac{\pi}{2} -\left(\alpha + \gamma\right) \right) =
		  r\sin\left(\frac{\pi}{2} + \left(\alpha + \gamma\right)\right) & \iff \\
		& \sin\left(\frac{\pi}{2} -\left(\alpha + \gamma\right) \right) =
		  \sin\left(\frac{\pi}{2} + \left(\alpha + \gamma\right)\right) & \iff \\
		& \sin\left(-\frac{\pi}{2} -\left(\alpha + \gamma\right) + \pi \right) =
		  \sin\left(\frac{\pi}{2} + \left(\alpha + \gamma\right)\right) & \iff \\
		& -\sin\left(-\frac{\pi}{2} -\left(\alpha + \gamma\right)\right) =
		  \sin\left(\frac{\pi}{2} + \left(\alpha + \gamma\right)\right) & \iff \\
		& \sin\left(\frac{\pi}{2} +\left(\alpha + \gamma\right)\right) =
		  \sin\left(\frac{\pi}{2} + \left(\alpha + \gamma\right)\right) \\
	\end{align*}

	Now, we need to prove that $B'_y > A'_y$ or $B'_y > C'_y$.
	\begin{align*}
		& r\sin\left(2\gamma - \alpha - \gamma + \frac{\pi}{2}\right) >
		  r\sin\left(-\alpha - \gamma + \frac{\pi}{2}\right) & \iff \\
		& \sin\left(2\gamma - \alpha - \gamma + \frac{\pi}{2}\right) >
		  \sin\left(-\alpha - \gamma + \frac{\pi}{2}\right) & \iff \\
		& \sin\left(\gamma - \alpha + \frac{\pi}{2}\right) >
		  \sin\left(-\alpha - \gamma + \frac{\pi}{2}\right)
	\end{align*}

	It is sufficient to prove that, for all $\alpha, \gamma \in \left(0, \pi\right)$, 
	$\sin\left(\gamma - \alpha + \frac{\pi}{2}\right) >
	 \sin\left(-\alpha - \gamma + \frac{\pi}{2}\right)$.
	Rearranging, we get
	$\sin\left(\gamma - \alpha + \frac{\pi}{2}\right) -
	 \sin\left(-\alpha - \gamma + \frac{\pi}{2}\right) > 0$.
	
	We can use the product-sum identity for sin to get the following:
	\begin{align*}
	& 2\cdot\cos\left(\frac{\gamma - \alpha + \frac{\pi}{2} - \alpha - \gamma + \frac{\pi}{2}}{2}\right) \cdot
		\sin\left(\frac{\gamma - \alpha + \frac{\pi}{2} + \alpha + \gamma - \frac{\pi}{2}}{2}\right) > 0 & \iff \\
	& 2\cdot\cos\left(\frac{-2\alpha + \pi}{2}\right) \cdot
		\sin\left(\frac{2\gamma}{2}\right) > 0 & \iff \\
	& \cos\left(-\alpha + \frac{\pi}{2}\right) \cdot
		\sin\left(\gamma\right) > 0
	\end{align*}

	Because $\gamma \in \left(0, \pi\right)$, $\sin\left(\gamma\right) > 0$, so it is sufficient to prove
	$\cos\left(\frac{\pi}{2} - \alpha\right) > 0 \iff \sin\left(\alpha\right) > 0$
	
	Because $\alpha \in \left(0, \pi\right)$, $\sin\left(\alpha\right) > 0$.

	Therefore, we have $\theta_{\HSview[B]} = -\alpha - \gamma + \frac{\pi}{2} - \theta = -\pi + \beta + \frac{\pi}{2}
	- \theta = \beta - \frac{\pi}{2} - \theta$.

	Finally, we observe $\HSview[A]$, we require that $a$ is parallel to the $x$-axis and that $A$ is above $a$ on the
	$y$-axis. Since a view is a rotation of the coordinate system, we will define the angle by which we need to rotate
	the coordinate system with $\theta_{\HSview[A]}$.
	
	Without loss of generality, let $\theta = 0$. If $\theta \neq 0$, we can
	simply add $-\theta$ to $\theta_{\HSview[A]}$ to undo its rotation.

	Let $A' = \HSview[A]\left(A\right), B' = \HSview[A]\left(B\right), C' = \HSview[A]\left(C\right)$.

	We require the following:
	$$
		A'_y > B'_y = C'_y \iff
		r\sin\left(\theta_{\HSview[A]}\right) >
		r\sin\left(2\gamma + \theta_{\HSview[B]}\right) =
		r\sin\left(2\gamma + 2\alpha + \theta_{\HSview[B]}\right)
	$$

	If we let $\theta_{\HSview[A]} = \beta - \gamma + \frac{\pi}{2}$, we get
	\begin{align*}
		& r\sin\left(2\gamma + \beta - \gamma + \frac{\pi}{2}\right) =
		  r\sin\left(2\gamma + 2\alpha + \beta - \gamma + \frac{\pi}{2}\right) & \iff \\
		& r\sin\left(\gamma + \beta + \frac{\pi}{2}\right) =
		  r\sin\left(\gamma + 2\alpha + \beta + \frac{\pi}{2}\right) & \iff \\
		& r\sin\left(\gamma + \beta + \frac{\pi}{2}\right) =
		  r\sin\left(\pi + \alpha + \frac{\pi}{2}\right) & \iff \\
		& r\sin\left(\gamma + \beta + \frac{\pi}{2}\right) =
		  r\sin\left(\pi + \pi - \beta - \gamma + \frac{\pi}{2}\right) & \iff \\
		& r\sin\left(\gamma + \beta + \frac{\pi}{2}\right) =
		  r\sin\left(2\pi - \beta - \gamma + \frac{\pi}{2}\right) & \iff \\
		& r\sin\left(\gamma + \beta + \frac{\pi}{2}\right) =
		  r\sin\left(- \beta - \gamma + \frac{\pi}{2}\right) & \iff \\
		& r\sin\left(\frac{\pi}{2} + \left(\beta + \gamma\right)\right) =
		  r\sin\left(\frac{\pi}{2} - \left(\beta + \gamma\right)\right) & \iff \\
		& r\sin\left(\frac{\pi}{2} + \left(\beta + \gamma\right)\right) =
		  r\sin\left(-\frac{\pi}{2} - \left(\beta + \gamma\right) + \pi\right) & \iff \\
		& r\sin\left(\frac{\pi}{2} + \left(\beta + \gamma\right)\right) =
		  -r\sin\left(-\left(\frac{\pi}{2} + \left(\beta + \gamma\right)\right)\right) & \iff \\
		& \sin\left(\frac{\pi}{2} + \left(\beta + \gamma\right)\right) =
		  \sin\left(\frac{\pi}{2} + \left(\beta + \gamma\right)\right)
	\end{align*}

	Now, we need to prove that $A'_y > B'_y$ or $A'_y > C'_y$.
	\begin{align*}
		& r\sin\left(\beta - \gamma + \frac{\pi}{2}\right) >
		  r\sin\left(2\gamma + \beta - \gamma + \frac{\pi}{2}\right) & \iff \\
		& r\sin\left(\beta - \gamma + \frac{\pi}{2}\right) >
		  r\sin\left(\gamma + \beta + \frac{\pi}{2}\right) & \iff \\
		& \sin\left(\beta - \gamma + \frac{\pi}{2}\right) >
		  \sin\left(\gamma + \beta + \frac{\pi}{2}\right)
	\end{align*}

	It is sufficient to prove that, for all $\beta, \gamma \in \left(0, \pi\right)$,
	$\sin\left(\beta - \gamma + \frac{\pi}{2}\right) >
	 \sin\left(\gamma + \beta + \frac{\pi}{2}\right)$. Rearranging, we get 
	$\sin\left(\beta - \gamma + \frac{\pi}{2}\right) -
	 \sin\left(\gamma + \beta + \frac{\pi}{2}\right) > 0$.

	We can use the product-sum identity for sin to get the following:
	\begin{align*}
	  & 2 \cdot \cos\left(\frac{\beta - \gamma + \frac{\pi}{2} + \gamma + \beta + \frac{\pi}{2}}{2}\right) \cdot
	    \sin\left(\frac{\beta - \gamma + \frac{\pi}{2} - \gamma - \beta - \frac{\pi}{2}}{2}\right) > 0 & \iff \\
	  & 2 \cdot \cos\left(\frac{2\beta + \pi}{2}\right) \cdot
	    \sin\left(\frac{-2\gamma}{2}\right) > 0 & \iff \\
	  & 2 \cdot \cos\left(\beta + \frac{\pi}{2}\right) \cdot
	    \sin\left(-\gamma\right) > 0 & \iff \\
	  & \cos\left(\beta + \frac{\pi}{2}\right) \cdot
	    \sin\left(-\gamma\right) > 0 & \iff \\
	  & -\cos\left(\beta + \frac{\pi}{2}\right) \cdot
	    \sin\left(\gamma\right) > 0 & \iff \\
	  & \cos\left(\beta + \frac{\pi}{2}\right) \cdot
	    \sin\left(\gamma\right) < 0
	\end{align*}
	
	Because $\gamma \in \left(0, \pi\right)$, $\sin\left(\gamma\right) > 0$, so it is sufficient to prove
	$\cos\left(\beta + \frac{\pi}{2}\right) < 0$

	Since 
	$\beta \in \left(0, \pi\right)$, $\beta + \frac{\pi}{2} \in \left(\frac{\pi}{2}, \frac{3\pi}{2}\right)$.
	Cosine is negative in that domain.


	Therefore, we have $\theta_{\HSview[A]} = \beta - \gamma + \frac{\pi}{2} - \theta$.

	Finally, we have the following:
	\begin{align*}
		\theta_{\HSview[A]} &= \beta - \gamma + \frac{\pi}{2} - \theta \\
		\theta_{\HSview[B]} &= \beta - \frac{\pi}{2} - \theta \\
		\theta_{\HSview[C]} &= \frac{3\pi}{2} - \gamma - \theta
	\end{align*}
	
	$\HSview[A]$ is a rotation of a point around the origin by $\theta_{\HSview[A]}$.\\
	$\HSview[B]$ is a rotation of a point around the origin by $\theta_{\HSview[B]}$.\\
	$\HSview[C]$ is a rotation of a point around the origin by $\theta_{\HSview[C]}$.

	Now, in the future, we do all of the calculations of further triangles on some side $v$ in view $\HSview$.
	After that, to get the actual position of the triangle in untransformed space, we apply $\HSview^{-1}$
	to the calculated triangle.
	
	\newpage

	\section{Base triangles}
	
	Suppose we are in some view transformed space $\HSview$.
	Let $v$ be a root side identifier part which matches $\HSview$ (i.e. $\HSview = \widetilde{A} \implies v = \HSindexpart{a}$, etc.).

	Because we are in $\HSview$, the bottom side of the root and any base triangle touches the circle at both of its vertices and is parallel to the $x$ axis.
	Therefore, the bottom points must be inverses of each other:
	\begin{equation}
		\HSLx < 0 < \HSRx = -\HSLx, \forall I \in \mathcal{I}_R \cup \mathcal{I}_{Bv} \label{base-left-right-identity}
	\end{equation}

	Since the bottom side touches the circle, and the right point is greater than zero, we conclude:
	\begin{equation}
		\HSRx = \sqrt{r^2 - \HSRy^2}, \forall I \in \mathcal{I}_R \cup \mathcal{I}_{Bv} \label{base-bottom-circle-identity}
	\end{equation}

	Because we concluded that the top point of a base triangle must touch its parent's bottom side, we have
	\begin{equation}
		\HSTy = \HSLy[P(I)] = \HSRy[P(I)], \forall I \in \mathcal{I}_{Bv} \label{base-top-bottom-identity}
	\end{equation}

	This also implies that, to get a triangle's height, we can subtract the y coordinate of its top point from the y coordinate of the top of the next base triangle's point.
	
	Let $I \in \mathcal{I}_{Bv}$. Assume coordinates of all of $\HStriangle[P(I)]$'s vertices are known. This is, in fact, true for the root triangle's coordinates. We will thus set up a recurrence relation.

	We know that any positive triangle's base side length is its right point's $x$ coordinate minus its left point's $x$ coordinate, so the side length of our known triangle is
	$$\HSRx[P(I)] - \HSLx[P(I)] \xlongequal{\eqref{base-left-right-identity}} 2\HSRx[P(I)]$$

	We know that any positive triangle's height is its top point's y coordinate minus its left/right point's y coordinate, so the height of our known triangle is:
	$$\HSTy[P(I)] - \HSRy[P(I)] \xlongequal{\eqref{base-top-bottom-identity}} \HSTy[P(I)] - \HSTy$$

	We know, because of \eqref{base-bottom-circle-identity} and \eqref{base-top-bottom-identity},
	$$\HSRx[P(I)] = +\sqrt{r^2 - \HSRy[P(I)]^2} = \sqrt{r^2 - \HSTy^2}$$
	
	We know that the ratio of a triangle's base to its height is constant between similar triangles, and we named that ratio $\HSk$.
	We are going to use this fact to calculate all of the next triangle's vertices' coordinates.

	\subsection{The y value of the next top point}

	We write:
	\begin{align*}
	  & \HSRx[P(I)] - \HSLx[P(I)] = \HSk\left(\HSTy[P(I)] - \HSRy[P(I)]\right)\\
		\implies & 2\sqrt{r^2 - \HSTy^2} = \HSk\HSTy[P(I)] - \HSk\HSTy\\
		\implies & 4r^2 - 4\HSTy^2 = \HSk^2\HSTy[P(I)]^2 - 2\HSk^2\HSTy[P(I)]\HSTy + \HSk^2\HSTy^2\\
		\implies & \left(k^2_{\HSview} + 4\right)\HSTy^2 + \left(-2k^2_{\HSview}\HSTy[P(I)]\right)\HSTy + \left(k^2_{\HSview}\HSTy[P(S)]^2 - 4r^2\right)\\
		\implies & \HSTy = \frac{
			2k^2_{\HSview} \pm \sqrt{
				4k^4_{\HSview}\HSTy[P(I)]^2 -
				4\left(
					k^4_{\HSview}\HSTy[P(I)]^2 -
					4r^2k^2_{\HSview} +
					4k^2_{\HSview}\HSTy[P(I)]^2 -
					16r^2
				\right)
			}
		}{2k^2_{\HSview} + 8}\\
		\implies & \HSTy = \frac{
			k^2_{\HSview}\HSTy[P(I)] \pm 2\sqrt{
				r^2k^2_{\HSview} -
				k^2_{\HSview}\HSTy[P(I)]^2 +
				4r^2
			}
		}{k^2_{\HSview} + 4}
	\end{align*}
	
	Since all of our top points have to be inside the circle---except the root triangle's top point, which may at most be at coordinate $(0, r)$---we set the following constraint:
	\begin{equation}
		\HSTy[P(I)] \in \left(-r, r\right]\label{base-circle-constraint}
	\end{equation}

	\begin{remark}
		The inclusion of $\HSTy[P(I)] = r$ is specially for when $P(I) = R$. For every other base index $I$, we will get $\HSTy \in \left(-r, r\right).$
		Notice also how $\eqref{base-circle-constraint}$ immediately implies that the discriminant is positive, since it has to be at least $4r^2$.
	\end{remark}

	Since the next base triangle has to be below the previous, our solution also has to satisfy the following:

	\begin{equation}
		\HSTy[P(I)] > \HSTy \label{base-height-requirement}
	\end{equation}


	\Needspace{60\baselineskip}
	\begin{proposition}
		The positive branch cannot satisfy both \eqref{base-circle-constraint} and \eqref{base-height-requirement} at once.
	\end{proposition}
	\begin{proof}
		Assume that the branch is positive and that \eqref{base-circle-constraint} and \eqref{base-height-requirement} both hold.
		\begin{align*}
			\eqref{base-height-requirement} \implies & \HSTy[P(I)] > \frac{k^2_{\HSview}\HSTy[P(I)] + 2\sqrt{k^2_{\HSview}r^2 - k^2_{\HSview}\HSTy[P(I)]^2 + 4r^2}}{k^2_{\HSview} + 4}\\
			\implies & k^2_{\HSview}\HSTy[P(I)] + 4\HSTy[P(I)] > k^2_{\HSview}\HSTy[P(I)] + 2\sqrt{k^2_{\HSview}r^2 - k^2_{\HSview}\HSTy[P(I)]^2 + 4r^2}\\
			\implies & 2\HSTy[P(I)] > \sqrt{k^2_{\HSview}r^2 - k^2_{\HSview}\HSTy[P(I)]^2 + 4r^2}
		\end{align*}
		We do case analysis on $\HSTy[P(I)]$.
		\begin{itemize}
			\item \textbf{Case 1: $-r < \HSTy[P(I)] \le 0$:}\\
				Then, $2\HSTy[P(I)] \leq 0$

				But, $\sqrt{k^2_{\HSview}r^2 - k^2_{\HSview}\HSTy[P(I)]^2 + 4r^2} > 0$

				$\contradiction$ The inequality cannot hold. 
			\item \textbf{Case 2: $r \ge \HSTy[P(I)] > 0$:}\\
				Since both sides of the inequality are positive, we can square them.
				\begin{align*}
					& 2\HSTy[P(I)] > \sqrt{k^2_{\HSview}r^2 - k^2_{\HSview}\HSTy[P(I)]^2 + 4r^2}\\
					\implies & 4\HSTy[P(I)]^2 > k^2_{\HSview}r^2 - k^2_{\HSview}\HSTy[P(I)]^2 + 4r^2\\
					\implies & \left(k^2_{\HSview} + 4\right)\HSTy[P(I)]^2 > \left(k^2_{\HSview} + 4\right)r^2\\
					\implies & \HSTy[P(I)]^2 > r^2\\
					\implies & \HSTy[P(I)] \in \left(r, \infty\right) &&\because \HSTy[P(I)] > 0
				\end{align*}

				$\contradiction$ If \eqref{base-height-requirement} holds, \eqref{base-circle-constraint} does not.
		\end{itemize}
	\end{proof}

	\newpage

	\begin{proposition}
		If the negative branch satisfies \eqref{base-circle-constraint}, it also satisfies \eqref{base-height-requirement}.
	\end{proposition}
	\begin{proof}
		Assume that the branch is negative and that \eqref{base-circle-constraint} holds.
		We need to prove
		\begin{align*}
			&\HSTy[P(I)] > \frac{k^2_{\HSview}\HSTy[P(I)] - 2\sqrt{k^2_{\HSview}r^2 - k^2_{\HSview}\HSTy[P(I)]^2 + 4r^2}}{k^2_{\HSview} + 4}\\
			\iff& k^2_{\HSview}\HSTy[P(I)] + 4\HSTy[P(I)] > k^2_{\HSview}\HSTy[P(I)] - 2\sqrt{k^2_{\HSview}r^2 - k^2_{\HSview}\HSTy[P(I)]^2 + 4r^2}\\
			\iff & 2\HSTy[P(I)] > -\sqrt{k^2_{\HSview}r^2 - k^2_{\HSview}\HSTy[P(I)]^2 + 4r^2}
		\end{align*}

		We do case analysis on $\HSTy[P(I)]$.
		\begin{itemize}
			\item \textbf{Case 1: $0 \le \HSTy[P(I)] \le r$:}\\
				$2\HSTy[P(I)] > -\sqrt{k^2_{\HSview}r^2 - k^2_{\HSview}\HSTy[P(I)]^2 + 4r^2}$
				trivially holds because LHS is nonnegative and RHS is negative.

			\item \textbf{Case 2: $-r < \HSTy[P(I)] < 0$:}\\
				\begin{align*}
					&\HSTy[P(I)] \in \left(-r, 0\right)\\
					\implies & \HSTy[P(I)]^2 < r^2\\
					\implies & \left(k^2_{\HSview} + 4\right)\HSTy[P(I)]^2 < \left(k^2_{\HSview} + 4\right)r^2\\
					\implies & 4\HSTy[P(I)]^2 < k^2_{\HSview}r^2 - k^2_{\HSview}\HSTy[P(I)]^2 + 4r^2\\
					\implies & -2\HSTy[P(I)] < \sqrt{k^2_{\HSview}r^2 - k^2_{\HSview}\HSTy[P(I)]^2 + 4r^2} && \because \HSTy[P(I)] < 0\\
					\implies&2\HSTy[P(I)] > -\sqrt{k^2_{\HSview}r^2 - k^2_{\HSview}\HSTy[P(I)]^2 + 4r^2}\\
				\end{align*}

		\end{itemize}
	\end{proof}
	

	Since this is a recurrence relation, our solution also has to satisfy the following:
	$$\HSTy \in \left(-r, r\right]$$
	In fact, we are going to prove it satisfies the following:
	\begin{equation}
		\HSTy \in \left(-r, r\right)\label{base-circle-constraint-induction-step}
	\end{equation}
	
	\newpage

	\begin{proposition}
		$\eqref{base-circle-constraint} \implies \eqref{base-circle-constraint-induction-step}$.
	\end{proposition}
	

	\begin{proof}
		We define a function $f : \left[-r, r\right] \to \mathbb{R}$ with the rule
		$$f(y) = \frac{
			k^2_{\HSview}y - 2\sqrt{
				r^2k^2_{\HSview} -
				k^2_{\HSview}y^2 +
				4r^2
			}
		}{k^2_{\HSview} + 4}$$

		We set its derivative equal to zero.

		\begin{align*}
			&\frac{1}{k^2_{\HSview} + 4}\left(k^2 - \frac{-2k^2_{\HSview}y}{\sqrt{r^2k^2_{\HSview} - k^2_{\HSview}y^2 + 4r^2}}\right) = 0\\
			\iff&\frac{y}{\sqrt{r^2k^2_{\HSview} - k^2_{\HSview}y^2 + 4r^2}} = -\frac{1}{2}
		\end{align*}

		Because of \eqref{base-circle-constraint}, we know that the denominator is positive, so y has to be negative for the expression to equal $-\frac{1}{2}$.

		We square both sides.
		
		\begin{align*}
			&\frac{y^2}{r^2k^2_{\HSview} - k^2_{\HSview}y^2 + 4r^2} = \frac{1}{4}\\
			\iff&r^2\left(k^2_{\HSview} + 4\right) - y^2\left(k^2_{\HSview} + 4\right) = 0 && \text{denominator is positive}
		\end{align*}

		Since $y$ is negative, the solution is $-r$, which is the only critical point of the function.

		Since $f$ is continuous on its domain, we can see that, because $f'(r) = \frac{2k^2_{\HSview}}{k^2_{\HSview} + 4} > 0$,
		the function is strictly increasing, and its minimum is at $-r$.

		So to get the minimum and maximum, we can just compute $f(-r) = -r$ and $f(r) = \frac{k^2_{\HSview} - 4}{k^2_{\HSview} + 4}r$.
		$\frac{k^2_{\HSview} - 4}{k^2_{\HSview} + 4} \in \left(-1, 1\right)$ for $k > 0$

		Therefore, $y \in \left(-r, r\right] \implies f(y) \in \left(-r, r\right)$

	\end{proof}
	\newpage
	
	Finally, we can conclude that, if a valid recurrence relation exists, it has to be:
	
	\begin{equation}
		\HSTy = \frac{
			k^2_{\HSview}\HSTy[P(I)] - 2\sqrt{
				r^2k^2_{\HSview} -
				k^2_{\HSview}\HSTy[P(I)]^2 +
				4r^2
			}
		}{k^2_{\HSview} + 4} \label{base-triangle-top-y-equation}
	\end{equation}

	\begin{conjecture} \label{base-valid-solution}
		$$\HSTx[R]^2 + \HSTy[R]^2 = r^2 \implies \HSTx[RvB]^2 + \HSTy[RvB]^2 < r^2$$
		$$\forall I \in \mathcal{I}_{Bv}, \HSTx^2 + \HSTy^2 < r^2 \implies \HSTx[I\doubleplus B]^2 + \HSTy[I\doubleplus B]^2 < r^2$$

		I.e. if the root triangle's top point is valid, every subsequent base triangle's top point has to stay inside the circle.
	\end{conjecture}

	\pagebreak

	\subsection{The rest of the triangle}

	Now that we have \eqref{base-triangle-top-y-equation}, we can use it to calculate everything else.

	First, let us assign \eqref{base-triangle-top-y-equation} to some function.
	Let $f : \left(-r, r\right] \to \left(-r, r\right)$ be defined by $$f(y) = \frac{k^2_{\HSview}y - 2\sqrt{r^2k^2_{\HSview} - k^2_{\HSview}y^2 + 4r^2}}{k^2_{\HSview} + 4}$$
	
	Then,
	$$\HSTy = f\left(\HSTy[P(I)]\right)$$
	$$\HSLy = \HSRy = \HSTy[I \doubleplus B] = f\left(\HSTy\right)
	\because \eqref{base-top-bottom-identity}$$
	$$\HSRx = \sqrt{r^2 - \HSRy^2}
	\because \eqref{base-bottom-circle-identity}$$
	$$\HSLx = -\HSRx \because \eqref{base-left-right-identity}$$

	Finally, for $\HSTx$, we know that the line passing through
	$\HST$ and $\HSL$ is of the angle $\HSlambda$.
	We also know that the angle passing through $\HSR$ and $\HST$ is of the angle $-\HSrho$, as it is mirrored.

	Then, we can set up an equation for $\HSTx$ using $\HSs[l] = \cot\left(\HSlambda\right)$ or $\HSs[r] = \cot\left(-\HSrho\right)$:
	$$\HSTx = \HSs[l]\left(\HSTy - \HSLy\right) + \HSLx = \HSs[r]\left(\HSTy - \HSRy\right) + \HSRx$$
	
	Finally, we can reverse the $T$, $L$, $R$ substitution to get the actual points:
	
	$\HSA = \begin{cases}
		\HST, & \HSview \text{ is } \HSview[A]\\
		\HSR, & \HSview \text{ is } \HSview[B]\\
		\HSL, & \HSview \text{ is } \HSview[C]\\
  \end{cases}$,\quad\quad
	$\HSB = \begin{cases}
		\HSL, & \HSview \text{ is } \HSview[A]\\
		\HST, & \HSview \text{ is } \HSview[B]\\
		\HSR, & \HSview \text{ is } \HSview[C]\\
  \end{cases}$,\quad\quad
	$\HSC = \begin{cases}
		\HSR, & \HSview \text{ is } \HSview[A]\\
		\HSL, & \HSview \text{ is } \HSview[B]\\
		\HST, & \HSview \text{ is } \HSview[C]\\
  \end{cases}$
	
	\pagebreak

	\section{Normal triangles}
	
	Suppose we are in some view transformed space $\HSview$.
	Let $v$ be a root side identifier part which matches $\HSview$ (i.e. $\HSview = \widetilde{A} \implies v = \HSindexpart{a}$, etc.).

	Because we are in $\HSview$, the bottom side of the normal triangle is parallel to the $x$ axis, so we can write:
	\begin{equation}
		\HSRy = \HSLy, \forall I \in \mathcal{I}_{Nvd}, d \in \set{\HSindexpart{l}, \HSindexpart{r}} \label{normal-left-right-identity}
	\end{equation}

	For some $I \in \mathcal{I}_{Nvd}$, where $d \in \set{\HSindexpart{l}, \HSindexpart{r}}$, we define the following:
	
	The base side sign constant:
	$$\HSdelta \coloneqq \begin{cases}-1,& d = \HSindexpart{l}\\1,& d = \HSindexpart{r}\end{cases} \implies \HSdelta^2 = 1$$

	The bottom side point that touches the circle:
	$$\HSPSI = \begin{cases}
		\HSL, & d = \HSindexpart{l}\\
		\HSR, & d = \HSindexpart{r}
	\end{cases}$$

	The bottom side point that does not touch the circle:
	$$\HSPHI = \begin{cases}
		\HSR, & d = \HSindexpart{l}\\
		\HSL, & d = \HSindexpart{r}
	\end{cases}$$

	Notice how $\eqref{normal-left-right-identity} \iff \HSPHIy = \HSPSIy$.

	Suppose we are on the left base side.
	We can make an equation for $\HSRy$ through $\HSO$. We know that the angle is supposed to be $\HSlambda$, so we write:
	$$\HSRx = \cot\left(\HSlambda\right)\left(\HSRy - \HSOy\right) + \HSOx$$

	Since we are on the left side, we have
	$$\HSPHIx = \HSs[l]\left(\HSPHIy - \HSOy\right) + \HSOx$$
	
	Suppose we are on the right base side.
	We can make an equation for $\HSLy$ through $\HSO$. We know that the angle is supposed to be $-\HSrho$, so we write:
	$$\HSLx = \cot\left(-\HSrho\right)\left(\HSLy - \HSOy\right) + \HSOx$$

	Since we are on the right side, we have
	$$\HSPHIx = \HSs[r]\left(\HSPHIy - \HSOy\right) + \HSOx$$

	We notice how this is equivalent to writing the same equation for both base sides:
	\begin{equation}
		\HSPHIx = \HSs\left(\HSPHIy - \HSOy\right) + \HSOx \label{normal-phi-x-identity}
	\end{equation}
	
	Suppose we are on the left base side. Then, the left point a normal triangle touches the circle.
	Notice how, because of the way the first base triangle's bottom points split the circular segment's arc
	(the base's next base child has inverse left and right points), the $x$ value of the left point is always negative.
	Analogously for the right base side, the $x$ coordinate of the right point is always positive.

	Keeping that in mind because $\HSPSI$ is on the circle, we can derive $\HSPSIx$ in terms of $\HSPSIy \xlongequal{\eqref{normal-left-right-identity}} \HSPHIy$:
	\begin{equation}
		\HSPSIx = \HSdelta \sqrt{r^2 - \HSPHIy^2} \label{normal-psi-x-identity}
	\end{equation}
	
	Notice that, regardless of the base side, the length of the bottom side of any triangle is $\HSRx - \HSLx$.
	Rewriting that in terms of $\HSPSI$ and $\HSPHI$, we have

	$$\begin{cases}
		\HSPHIx - \HSPSIx, & d = \HSindexpart{l}\\\HSPSIx - \HSPHIx, & d = \HSindexpart{r}
	\end{cases}$$

	or equivalently, $\delta_d\left(\HSPSIx - \HSPHIx\right)$

	Notice that, regardless of the base side, the height from the top point of any triangle is $\HSTy - \HSPHIy$ or equivalently,
	$\HSOy - \HSPHIy$

	We can now use the same method we used for base triangles to set up a quadratic equation for $\HSPHIy$.
	\newpage
	\subsection{The y value of the bottom side point touching a triangle}
	
	We are going to define helper constants that will appear in the derivation:
	$$\HSsigma \coloneqq \HSdelta \HSs - \HSk$$
	$$\HSkappa \coloneqq \HSdelta \HSOx - \HSsigma \HSOy$$

	The second constant is for brevity and will be expanded again after we get the solution.

	We write:
	\begin{align*}
		         & \HSdelta \left( \HSPSIx - \HSPHIx \right) = \HSk \left( \HSOy - \HSPHIy \right)\\
		\implies & \sqrt{r^2 - \HSPHIy^2} - \HSdelta \HSs \HSPHIy + \HSdelta \HSs \HSOy - \HSdelta \HSOx = \HSk \HSOy - \HSk \HSPHIy\\
		\implies & \sqrt{r^2 - \HSPHIy^2} = \left(\HSdelta \HSs - \HSk\right)\HSPHIy + \HSdelta\HSOx - \left(\HSdelta \HSs - \HSk\right)\HSOy\\
		\implies & \sqrt{r^2 - \HSPHIy^2} = \HSsigma \HSPHIy + \left(\HSdelta\HSOx - \HSsigma\HSOy\right)\\
		\implies & r^2 - \HSPHIy^2 = \HSsigma^2 \HSPHIy^2 + 2 \HSsigma\left(\HSdelta\HSOx - \HSsigma\HSOy\right)\HSPHIy + \left(\HSdelta\HSOx - \HSsigma\HSOy\right)^2\\
		\implies & \left(\HSsigma^2 + 1\right)\HSPHIy^2 + \left(2\HSsigma\HSkappa\right)\HSPHIy + \left(\HSkappa^2 - r^2\right) = 0\\
		\implies &\HSPHIy = \frac{-2\HSsigma\HSkappa \pm \sqrt{4\HSsigma^2\HSkappa^2 - 4\left(\HSsigma^2 + 1\right)\left(\HSkappa^2 - r^2\right)}}{2\HSsigma^2 + 2}\\
		\implies &\HSPHIy = \frac{-\HSsigma\HSkappa \pm \sqrt{r^2\left(\HSsigma^2 + 1\right) - \HSkappa^2}}{\HSsigma^2 + 1}\\
		\implies &\HSPHIy = \frac{\HSsigma^2\HSOy - \HSdelta\HSsigma\HSOx \pm \sqrt{r^2\left(\HSsigma^2 + 1\right) - \left(\HSdelta \HSOx - \HSsigma \HSOy\right)^2}}{\HSsigma^2 + 1}\\
	\end{align*}

	Since the pseudo-sector origin point $\HSO$ has to be inside the circle, we set the following constraint:
	\begin{equation}
		\HSOx^2 + \HSOy^2 < r^2 \label{normal-circle-constraint}
	\end{equation}

	Since the normal triangle we are constructing has to be below the previous, our solution also has to satisfy the following:
	\begin{equation}
		\HSPHIy < \HSOy \label{normal-height-requirement}
	\end{equation}

	\newpage

	First, we have to prove that any solution is going to be real.
	\begin{proposition} \label{normal-real-solution-exists}
		$\eqref{normal-circle-constraint}$ implies that the expression under the square root of the solution is positive.
	\end{proposition}

	\begin{proof}
		First, let's prove that the expression cannot equal 0.
		Assume
		\begin{align*}
			r^2\left(\HSsigma^2 + 1\right) - \left(\HSdelta \HSOx - \HSsigma \HSOy\right)^2 = 0
			\implies r^2\left(\HSsigma^2 + 1\right) = \left(\HSdelta \HSOx - \HSsigma \HSOy\right)^2
		\end{align*}
		Now, we write
		{
			\setlength{\jot}{0.75em}
			\begin{align*}
			         	 & \HSOx^2 + \HSOy^2 < r^2\\
				\implies & \left(\HSOx^2 + \HSOy^2\right)\left(\HSsigma^2 + 1\right) < r^2\left(\HSsigma^2 + 1\right)\\
				\implies & \left(\HSOx^2 + \HSOy^2\right)\left(\HSsigma^2 + 1\right) < \left(\HSdelta \HSOx - \HSsigma \HSOy\right)^2\\
				\implies & \HSsigma^2\HSOx^2 + \HSsigma^2\HSOy^2 + \HSOx^2 + \HSOy^2 < \left(\HSdelta \HSOx - \HSsigma \HSOy\right)^2\\
				\implies & \HSsigma^2\HSOx^2 + \HSsigma^2\HSOy^2 + \HSOx^2 + \HSOy^2 < \HSOx^2 - 2\HSdelta\HSsigma\HSOx\HSOy  + \HSsigma^2 \HSOy^2\\
				\implies & \HSsigma^2\HSOx^2 + 2\HSdelta\HSsigma\HSOx\HSOy + \HSOy^2 < 0\\
				\implies & \left(\HSdelta\HSsigma\HSOx + \HSOy\right)^2 < 0 \contradiction
			\end{align*}
		}

		Therefore, the expression can never be exactly $0$.
		
		To prove it cannot be negative, we will use the intermediate value theorem.
		With $\HSdelta \in \set{-1, 1}$ and $r > 0$ constant,
		we define a domain
		$$D = \set{(x, y, \sigma) \in \mathbb{R}^3 \given x^2 + y^2 < r^2}$$
		and we define
		$$f : D \to \mathbb{R},\quad f(x, y, \sigma) = r^2\left(\sigma^2 + 1\right) - \left(\HSdelta x - \sigma y\right)^2$$
		Notice how we have proven that this function cannot equal $0$.

		Since $D$ is path-connected and $f$ is continuous, we evaluate $f\left(0, 0, 0\right) = r^2 > 0$.
		Suppose that there exists some point
		$\left(\overline{x}, \overline{y}, \overline{\sigma}\right)$ where
		$f\left(\overline{x}, \overline{y}, \overline{\sigma}\right) < 0$.

		Since $D$ is path-connected, for every path function $p : \left[0, 1\right] \to D$ with
		$p\left(0\right) = \left(0, 0, 0\right)$ and $p\left(1\right) = \left(\overline{x}, \overline{y}, \overline{\sigma}\right)$,
		$f \circ p : \left[0, 1\right] \to \mathbb{R}$ is continuous. Then, because
		$\left(f \circ p\right)\left(0\right) > 0$ and $\left(f \circ p\right)\left(1\right) < 0$,
		we can apply the intermediate value theorem to conclude
		that there must exist some $\xi \in \left(0, 1\right)$ such that $\left(f \circ p\right)\left(\xi\right) = 0$.

		We have already proven that such a point cannot exist. Therefore, neither can $\left(\overline{x}, \overline{y}, \overline{\sigma}\right)$.
		Since the function is never $0$, and there cannot exist a point where it is negative, it must be positive on its domain.
		Therefore, the expression under the square root equivalent to the function is also positive.
	\end{proof}

	\newpage
	
	\begin{proposition}
		The positive branch cannot satisfy both \eqref{normal-circle-constraint} and \eqref{normal-height-requirement} at the same time.
	\end{proposition}

	\begin{proof}
		Assume that the positive branch is the solution and that \eqref{normal-circle-constraint} and \eqref{normal-height-requirement} both hold.

		Let $S = \sqrt{r^2\left(\HSsigma^2 + 1\right) - \left(\HSdelta \HSOx - \HSsigma \HSOy\right)^2}$

		Then,
		\begin{align*}
			&\eqref{normal-height-requirement} \iff \HSOy > \frac{\HSsigma^2\HSOy - \HSdelta\HSsigma\HSOx + S}{\HSsigma^2 + 1}\\
			\iff& \HSOy\HSsigma^2 + \HSOy > \HSsigma^2\HSOy - \HSdelta\HSsigma\HSOx + S\\
			\iff& \HSOy + \HSdelta\HSsigma\HSOx > S\\
			\iff& \HSOy + \HSdelta\HSsigma\HSOx > \sqrt{r^2\left(\HSsigma^2 + 1\right) - \left(\HSdelta \HSOx - \HSsigma \HSOy\right)^2}
		\end{align*}
		Since \eqref{normal-circle-constraint} holds, the RHS is positive by \ref{normal-real-solution-exists}.

		We do case analysis on $\HSOy + \HSdelta\HSsigma\HSOx$.
		\begin{itemize}
			\item \textbf{Case 1: $\HSOy + \HSdelta\HSsigma\HSOx \le 0$:}\\

			$\contradiction$ The inequality cannot hold.
			\item \textbf{Case 2: $\HSOy + \HSdelta\HSsigma\HSOx > 0$:}\\

				Since LHS and RHS are both positive, we can square both sides and preserve the inequality.
				\begin{align*}
					&\HSOy + \HSdelta\HSsigma\HSOx > \sqrt{r^2\left(\HSsigma^2 + 1\right) - \left(\HSdelta \HSOx - \HSsigma \HSOy\right)^2}\\[1em]
					\iff & \HSOy^2 + 2\HSdelta\HSsigma\HSOy\HSOx + \HSsigma^2\HSOx^2 >{}\\
					{}>{}& r^2\left(\HSsigma^2 + 1\right) - \HSOx^2 + 2\HSdelta\HSsigma\HSOy\HSOx - \HSsigma^2\HSOy^2\\[1em]
					\iff & \HSOx^2 + \HSOy^2 + \HSsigma^2\HSOx^2 + \HSsigma^2\HSOy^2 > r^2\left(\HSsigma^2 + 1\right)\\[1em]
					\iff & \left(\HSOx^2 + \HSOy^2\right)\left(\HSsigma^2 + 1\right) > r^2\left(\HSsigma^2 + 1\right)\\[1em]
					\iff & \HSOx^2 + \HSOy^2 > r^2
				\end{align*}
				$\contradiction$ This is incompatible with \eqref{normal-circle-constraint}, which we assumed holds.
		\end{itemize}

	\end{proof}

	\newpage
	
	\begin{proposition}
		If the negative branch satisfies \eqref{normal-circle-constraint}, it satisfies \eqref{normal-height-requirement}.
	\end{proposition}

	\begin{proof}
		Assume that the negative branch is the solution and that \eqref{normal-circle-constraint} holds.

		Let $S = -\sqrt{r^2\left(\HSsigma^2 + 1\right) - \left(\HSdelta \HSOx - \HSsigma \HSOy\right)^2}$

		Then, we need to prove
		\begin{align*}
			&\eqref{normal-height-requirement} \iff \HSOy > \frac{\HSsigma^2\HSOy - \HSdelta\HSsigma\HSOx + S}{\HSsigma^2 + 1}\\
			\iff& \HSOy\HSsigma^2 + \HSOy > \HSsigma^2\HSOy - \HSdelta\HSsigma\HSOx + S\\
			\iff& \HSOy + \HSdelta\HSsigma\HSOx > S\\
			\iff& \HSOy + \HSdelta\HSsigma\HSOx > -\sqrt{r^2\left(\HSsigma^2 + 1\right) - \left(\HSdelta \HSOx - \HSsigma \HSOy\right)^2}\\
			\iff& -\HSOy - \HSdelta\HSsigma\HSOx < \sqrt{r^2\left(\HSsigma^2 + 1\right) - \left(\HSdelta \HSOx - \HSsigma \HSOy\right)^2}
		\end{align*}
		Since \eqref{normal-circle-constraint} holds, the RHS is positive by \ref{normal-real-solution-exists}.

		We do case analysis on $-\HSOy - \HSdelta\HSsigma\HSOx$.
		\begin{itemize}
			\item \textbf{Case 1: $-\HSOy - \HSdelta\HSsigma\HSOx \le 0$:}\\
			
			The inequality trivially holds.
			\item \textbf{Case 2: $-\HSOy - \HSdelta\HSsigma\HSOx > 0$:}\\
				\begin{align*}
					&\HSOy^2 + \HSOy^2 < r^2\\[1em]
					\implies & \left(\HSOy^2 + \HSOy^2\right)\left(\HSsigma^2 + 1\right) < r^2\left(\HSsigma^2 + 1\right)\\[1em]
					\implies & \HSOx^2 + \HSOy^2 + \HSsigma^2\HSOx^2 + \HSsigma^2\HSOy^2 < r^2\left(\HSsigma^2 + 1\right)\\[1em]
					\implies & \HSOy^2 + 2\HSdelta\HSsigma\HSOy\HSOx + \HSsigma^2\HSOx^2 <{}\\
					{}<{}& r^2\left(\HSsigma^2 + 1\right) - \HSOx^2 + 2\HSdelta\HSsigma\HSOy\HSOx - \HSsigma^2\HSOy^2\\[1em]
					\implies&\left(\HSOy + \HSdelta\HSsigma\HSOx\right)^2 < r^2\left(\HSsigma^2 + 1\right) - \left(\HSdelta \HSOx - \HSsigma \HSOy\right)^2\\[1em]
					\implies& -\HSOy - \HSdelta\HSsigma\HSOx < \sqrt{r^2\left(\HSsigma^2 + 1\right) - \left(\HSdelta \HSOx - \HSsigma \HSOy\right)^2}
				\end{align*}

		\end{itemize}

	\end{proof}

	Finally, we conclude that, if a valid recurrence relation exists, it has to be:
  \begin{equation}
  	\HSPHIy = \frac{\HSsigma^2\HSOy - \HSdelta\HSsigma\HSOx - \sqrt{r^2\left(\HSsigma^2 + 1\right) - \left(\HSdelta \HSOx - \HSsigma \HSOy\right)^2}}{\HSsigma^2 + 1} \label{normal-triangle-phi-y-equation}
  \end{equation}
	
	\begin{conjecture} \label{normal-valid-solution}
		Assuming that $\ref{base-valid-solution}$ holds,
		for some base index $I_B$ and all normal indices $I$, where $I$ is a descendent of $I_B$,
		$$\HSOx^2 + \HSOy^2 < r^2 \implies \HSPHIx^2 + \HSPHIy^2 < r^2$$
		$$\HSOx^2 + \HSOy^2 < r^2 \implies \HSTx^2 + \HSTy^2 < r^2$$

		I.e. if the base triangle $\HStriangle[I_B]$ that $\HStriangle$ is a descendent of is valid,
		then any of its normal descendents is valid as well.
	\end{conjecture}

	In combination, \ref{base-valid-solution} and \ref{normal-valid-solution} imply that, if the root triangle is valid (non-degenerate), then all of its descendents are. Since we can always find a descendent using the base and normal formulas, there are an infinite amount of triangles that can fit into the circle for any choice of root triangle.

	The only `part' of these conjectures that I have successfully proven is that the base triangle's top vertex's $y$ coordinate remains
	in the range $\left(-r, r\right]$.
	
	\newpage


	\subsection{The rest of the triangle}
	Now that we have \eqref{normal-triangle-phi-y-equation}, we can use it to calculate everything else.

	First, let us assign \eqref{normal-triangle-phi-y-equation} to some function.
	Let $f : \left(-r, r\right) \times \left(-r, r\right) \to \left(-r, r\right)$ be defined by
  $$
  f(x, y) =
  \frac{\HSsigma^2y - \HSdelta\HSsigma x - \sqrt{r^2\left(\HSsigma^2 + 1\right) - \left(\HSdelta x - \HSsigma y\right)^2}}{\HSsigma^2 + 1}
  $$

	Then,
	$$\HSPHIy = f\left(\HSOx, \HSOy\right)$$
	$$\HSPHIx = \HSs\left(\HSPHIy, - \HSOy\right) + \HSOx \because \eqref{normal-phi-x-identity}$$
	$$\HSPSIy = \HSPHIy \because \eqref{normal-left-right-identity}$$
	$$\HSPSIx = \HSdelta \sqrt{r^2 - \HSPHIy^2} \because \eqref{normal-psi-x-identity}$$

	We can reverse our $\Phi$ and $\Psi$ substitutions:

	$$\HSL = \begin{cases}
		\HSPSI, & d = \HSindexpart{l}\\
		\HSPHI, & d = \HSindexpart{r}
	\end{cases}$$

	$$\HSR = \begin{cases}
		\HSPHI, & d = \HSindexpart{l}\\
		\HSPSI, & d = \HSindexpart{r}
	\end{cases}$$

	$$\HSTy = \HSOy$$

	For $\HSTx$, we know that the line passing through
	$\HST$ and $\HSL$ is of the angle $\HSlambda$.
	We also know that the angle passing through $\HSR$ and $\HST$ is of the angle $-\HSrho$, as it is mirrored.

	Then, we can set up an equation for $\HSTx$ using $\HSs[l] = \cot\left(\HSlambda\right)$ or $\HSs[r] = \cot\left(-\HSrho\right)$:
	$$\HSTx = \HSs[l]\left(\HSTy - \HSLy\right) + \HSLx = \HSs[r]\left(\HSTy - \HSRy\right) + \HSRx$$
	
	Finally, we can reverse the $T$, $L$, $R$ substitution to get the actual points:
	
	$\HSA = \begin{cases}
		\HST, & \HSview \text{ is } \HSview[A]\\
		\HSR, & \HSview \text{ is } \HSview[B]\\
		\HSL, & \HSview \text{ is } \HSview[C]\\
  \end{cases}$,\quad\quad
	$\HSB = \begin{cases}
		\HSL, & \HSview \text{ is } \HSview[A]\\
		\HST, & \HSview \text{ is } \HSview[B]\\
		\HSR, & \HSview \text{ is } \HSview[C]\\
  \end{cases}$,\quad\quad
	$\HSC = \begin{cases}
		\HSR, & \HSview \text{ is } \HSview[A]\\
		\HSL, & \HSview \text{ is } \HSview[B]\\
		\HST, & \HSview \text{ is } \HSview[C]\\
  \end{cases}$

	\newpage
	\section{Algorithm}
	This section is more informal than the previous ones and simply describes how the program bundled with this document
	functions using all of the information described above.
	
	The following describes the areas inside the circle that are not necessarily covered by triangles.
	
	It is easy to conclude that, in some view transformed space, if we
	take some normal triangle on the left side, and reflect it about its right side,
	the space left blank by creating that triangle will be perfectly filled.
	This is because the space left blank is a triangle, it shares one of the sides, and it can easily be shown
	that it shares all angles with the existing triangle. So, by side-angle-side, it has to be congruent.

	We will call such a triangle a \textbf{negative triangle}. Both because it is in the negative space
	left by adding a positive triangle into a pseudo-sector, or because it is flipped upside down from a
	positive triangle in the view transformation.

	So, in some view transformed space, we have the following primitive shapes: \begin{itemize}
		\item positive triangles
		\item negative triangles
		\item circular segments
		\item circular pseudo-sectors
	\end{itemize}

	We define a \textbf{shape collection} in to be one of:\begin{itemize}
		\item root shape collection -   a triangle and three circular segments corresponding to each of the triangle's sides,
		\item base shape collection -   a positive triangle, a circular segment corresponding to its bottom
		\item base side shape collection - one pseudo-segment
		\item normal shape collection - two triangles, one positive, one negative, and two pseudo-sectors, one horizontal, one vertical
	\end{itemize}

	Shape collections are simply a way to hold primitives that the algorithm may want to draw.

	The root triangle node holds a root shape collection labeled ``untransformed''.

	A root side node holds a base shape collection labeled ``transformed''. This node also holds the side-to-height ratio.

	A base triangle node holds a base shape collection labeled ``transformed'' and another labeled ``untransformed''.

	A base side node holds a base side shape collection labaled ``transformed'' and another labeled ``untransformed''. This node also holds the left/right slope.
	
	A normal triangle node holds a normal shape collection labeled ``transformed'' and another labeled ``untransformed''.

	The main program is split into multiple stages.

	The program is split into two stages.

	The first stage involves creating a \textbf{triangle tree} object using a radius $r$, and angles
	$\alpha$, $\beta$, $\theta$, described above.

	Then, we calculate the root, and recursively calculate its children. Creating a node populates ``transformed'' (if applicable) and ``untransformed'' (if applicable) with the proper shape collections.
	
	The root side nodes contain a
	view transformed root triangle of the respective side. Effectively, in a view transformed space, the root triangle becomes the `0th' base triangle, hence why the root side node itself has a base triangle shape collection.

	The tree in the algorithm acts identically to the one described in the sections above.

	The idea is that the transformed shape collections serve as a way to create new nodes, with their own
	shape collections derived from the previous ones. This is the \textbf{calculation} stage, where the algorithm recursively creates new nodes from existing nodes, using the
	equations we have described in previous sections.

	At this point, our triangle tree now has all of the primitive data required to draw the shape we want to draw.

	The next stage is the \textbf{drawing} stage, where an interface called a \texttt{Drawer} may do whatever it wants
	with the information passed from a triangle tree.

	A wrapper around a calculation and drawing sequence called a \texttt{Processor} has also been created in the example code, and used is to draw the figures in this document.

\end{document}
